\documentclass[10pt]{article}
\usepackage{template}
\usepackage{microtype,amsmath}
\setlength{\emergencystretch}{2em}
\SetKeys{date=2026-11-22,vol={2026-W47},title={每周重大新闻综述},author={ChatGPT 编辑部},note={\small 观察期：2026年11月15日至21日。截稿：北京时间11月22日上午。事实、各方主张、推论与未确认事项分开陈述。}}
\renewcommand{\newsitem}[3]{\par\medskip{\large\bfseries\color{newsblue}#1\par}{\footnotesize\color{newsgray}#2\par}\smallskip #3\par\medskip{\color{newsgray}\rule{\linewidth}{0.25pt}}\par}
\begin{document}
\begin{documents}
\newsitem{COP31：电气化目标为何无法代替化石能源转型}{深度综述｜11月19—22日｜路透社、卫报、AP}{
在土耳其安塔利亚举行的第三十一届联合国气候变化大会，原定于11月20日结束，却因为化石能源与气候资金问题持续延会。21日的草案谈判仍存在实质分歧；22日凌晨大会主席表示已提出可能获得共识的新文本，但在本刊截稿时尚不能将其视为正式通过的最终决议。会议程序中的“接近一致”与各国真正承担新义务之间仍有很长距离。\cite{cop_reuters,cop_ap}

争议首先来自一个看似积极的数字：到2035年把电力在全球终端能源消费中的占比提高至35\%。电气化可以将燃油汽车、燃气供暖和部分工业过程转向电力，但这并不自动减少排放。如果新增电力主要来自煤炭与天然气，发电环节的排放可能抵消终端使用效率提高带来的收益。评价该目标，必须同时观察发电结构、全生命周期排放和替代化石燃料的实际速度，而不是把用电占比直接解释为清洁能源占比。\cite{cop_power}

对出口油气的国家而言，明确的淘汰时间表可能触及财政收入、就业和既有基础设施的价值；对低海拔岛国及易受极端天气影响的国家而言，继续拖延则意味着更高的灾害与适应成本。欧盟与若干脆弱国家要求保留明确的化石燃料转型措辞，沙特阿拉伯及部分产油国则反对过于强硬的表述。双方争论的不只是修辞，而是未来投资信号、国际资金流向以及是否存在可衡量的减排义务。\cite{cop_reuters,cop_guardian}

气候融资是另一条断裂线。发展中国家要求富裕国家提供更明确的适应资金安排，因为海堤、预警系统、抗旱基础设施和医疗体系建设的收益主要体现为未来损失减少，短期内不一定产生可用于偿债的现金流。草案中曾出现的1200亿美元适应资金目标后来被删除，引起反对。融资争议涉及历史排放责任、当前支付能力、受援资金的使用透明度，以及贷款与赠款应如何组合；任何一个环节都无法用一句“增加援助”解决。\cite{cop_reuters,cop_ap}

大会还围绕碳捕集与移除技术产生争论。这些技术可能在水泥、钢铁和难以彻底消除的残余排放中发挥作用，但不能因此推断它们可以无成本、无限规模地替代直接减排。支持者强调技术组合的必要性，批评者担心模糊的未来移除承诺会为当下持续使用化石燃料提供理由。应当把技术的试验规模、实际永久封存量、能耗和成本与政治承诺分开核查。

人工智能基础设施也进入气候外交议程。数据中心的电力需求与供电可靠性要求，使电气化不再只是交通和家庭供暖的问题。即使模型单位推理能耗降低，使用量增加仍可能推动总负荷上升；电网扩容又受制于变压器、输电线路、许可和土地。所谓“绿色人工智能”若没有明确的用电时间、地点、边际电源与测量口径，就很容易退化成宣传标签。

本刊认为，衡量本次大会不应只看最终文件是否出现“化石燃料”四个字，也要看它是否形成可追踪的行动路径、明确资金承诺与定期复核机制。反过来，不能因为一份文本未能满足最强要求，就认定所有国家的实际减排政策同时归零。多边协议是协调工具，不是排放变化本身。未来应逐项追踪正式决议、各国提交的国家自主贡献、资金拨付和能源投资。截稿时协议尚在变动，最终条文必须以会后公布的正式文本为准。

}

\newsitem{乌克兰战争：打击升级与谈判议程同时扩张}{深度综述｜11月17—21日｜Reuters、AP、The Guardian}{
11月17日，乌克兰及其欧洲支持者要求尽快交付防空装备；20日，俄罗斯对乌克兰城市和能源、交通设施发动大规模导弹与无人机袭击。乌方官员报告基辅、利沃夫等地出现死伤和基础设施受损，并将攻击与克里米亚大桥遭袭后的报复行动联系起来。伤亡数字是特定时点的官方通报，后续调查可能调整；具体武器数量、打击目标和各方责任仍应以独立取证补充核实。\cite{ukr_reuters,ukr_guardian}

同一周，围绕战争结束后的欧洲安全架构，外交接触并未停止。美方提出讨论战后安全安排，俄罗斯将其视为重新提出部分长期安全诉求的机会。俄方反对乌克兰加入北约及欧洲军队在战后驻乌；乌克兰与欧洲国家担心，在停火和安全保证尚未明确前就谈判大规模调整北约部署，会使遭受侵略的一方失去谈判筹码。美方则强调希望欧洲参与并讨论类似北约的安全保证。各方所说的“安全”并不是同一个可直接相加的指标。\cite{ukr_ap}

谈判的第一层问题是停止敌对行动：何时停火、如何界定接触线、怎样监督违约、是否交换战俘和恢复基础服务。第二层是长期安全保证：若战争再次爆发，谁负有何种援助义务，是否部署部队或提供防空能力。第三层才涉及领土、制裁与国际秩序。将三层议题捆绑，可能使任何单一分歧阻断人道安排；完全拆开，又可能让临时停火缺乏持续动力。不同国家选择何种谈判顺序，反映了对违约风险的不同判断。

基础设施袭击直接影响平民生活。冬季供暖依赖电力和燃料供应，城市供水又依赖泵站与处理设施，铁路既是民用交通系统也是重要后勤网络。停电和供水中断可能造成的次生损失，不能仅由当天的死亡人数衡量。战争报道也必须区分军事设施、军民两用基础设施与纯民用目标，不能仅凭交战一方的描述直接作出法律结论；国际人道法关于区分、比例和预防措施的要求，需要针对具体袭击调查。

与此同时，乌克兰对克里米亚大桥等俄罗斯控制的交通设施采取打击，俄罗斯强调其军事与政治后果。双方均试图削弱对方的后勤与战略意志。这里应区分目标的实际军事用途、袭击手段和附带损害，不能仅凭某设施具有民用功能或军事功能之一，就推导所有攻击当然合法或非法。

欧洲各国的政策约束也并不一致。邻近冲突地区的国家更重视领土防御和威慑，其他国家还要平衡财政、能源和国内选举。美国在调停中具有重要军援与制裁筹码，但不能代替乌克兰决定其主权和安全安排。俄罗斯的长期目标与战场承受能力，亦不应只根据单次声明判断。

因此，未来一周真正值得观察的，不是“双方同意继续接触”这样的标题，而是是否出现有日期、地点、参与方和具体议程的正式会谈；是否建立可核查的人道安排；以及空袭强度是否随外交接触而下降。若战场行动与谈判持续同步升级，和平进程的可信度仍会受到削弱。

}

\newsitem{自主武器国际规则：技术进步快于共识机制}{深度综述｜11月21日｜Reuters、The Guardian}{
在日内瓦举行的《特定常规武器公约》框架内谈判，11月21日未能就致命自主武器形成有法律约束力的明确规则。与会国家同意继续制定可能的规范，目标时间延至2028年，但新安排是否具有法律约束力仍未确定。报道显示约百个国家支持较快制定约束性条约，美国、俄罗斯、以色列等主要军事技术国家反对相关方案。把结果称为“谈判彻底终止”并不准确：程序仍在继续，但最有争议的义务尚未达成一致。\cite{arms_reuters,arms_guardian}

致命自主武器并非简单等于无人机。关键问题是系统在多大程度上自行选择、识别和攻击目标，人类操作者是否能理解系统判断并及时中止。自动驾驶、目标跟踪、遥控武器与自主选择攻击目标之间存在连续谱；不同国家可能使用不同定义。这种定义分歧直接影响条约的适用范围，也决定软件升级或改变操作流程后，原有认证是否仍然有效。

支持强制性规则的国家和人权组织强调，机器学习系统在复杂战场中可能误认目标，训练数据也难以覆盖现实环境的全部变化。平民与战斗员的区分、攻击的比例原则以及责任追究，不能只由一个置信度阈值代替。若系统输出错误，究竟由指挥官、开发者、采购部门还是操作人员负责，必须在部署之前建立清晰制度。

反对新条约的国家则可能认为现有国际人道法已经适用于任何武器，另设规则会限制合法军事创新；有的国家担心定义过宽，把用于防御来袭导弹的自动拦截系统也纳入限制。这些顾虑值得具体讨论，但并不自动证明现有法律足以解决所有新风险。既有法律的原则能否转化为可操作的测试、记录和问责程序，仍是需要验证的独立问题。

《特定常规武器公约》依赖协商一致，少数反对方即可阻碍有约束力的成果。这使大国更容易保护其技术空间，却也促使一些国家考虑在其他国际论坛推进条约。历史上某些军控安排确曾在非全体大国参与的框架下先行建立规范，但其执行效果与军事大国是否加入密切相关。

技术上，测试自主武器不能只统计平均识别准确率。应检验罕见场景、传感器故障、对抗欺骗、模型更新后的行为漂移，以及通信中断时系统是否进入安全状态。尤其需要区分系统在实验室完成某种任务，与其在真实环境中符合国际人道法要求这两种不同结论。

本刊判断，未来规则至少需要讨论可审计的使用记录、明确的人类责任链、禁止或限制无法可靠区分目标的应用，以及对部署条件的持续审查。究竟采用普遍禁止、分级限制还是技术标准，仍属政治和法律选择，不能由工程师单独决定。

}

\newsitem{航天观察：旅行者一号达到单程光信号一天的距离}{重点报道｜NASA｜11月19日}{
美国国家航空航天局11月19日公布，旅行者一号已到达距离地球约一个光日的位置。这是距离尺度的里程碑，而不是飞船突然进入某个新物理区域。所谓一个光日，是光在真空中传播24小时的距离，约为259亿公里；在这一距离上，地球向飞船发送命令、飞船再回传结果，往返信号延迟约48小时。\cite{nasa_voyager}

这对任务运行提出了非常具体的要求。地面团队无法像控制近地卫星一样快速交互排错，每次指令必须预先考虑飞船的状态、能量预算和潜在故障。旅行者一号1977年发射，早已进入星际空间，所携带的计算机、无线电与电源系统仍然受当时工程技术的限制。随着放射性同位素电源输出逐年下降，团队必须权衡科学仪器与维持飞船基本运行的能耗。

里程碑也提醒读者，距离和时间在深空通信中不可分割。地面看到的遥测信息对应的是飞船一天之前的状态，通信频率、天线增益和信噪比共同决定可获得的数据量。距离增加会使接收信号更加微弱，不能仅凭一条成功的载波确认就推断所有科学仪器运行正常。任务状态应以NASA后续工程公告为准。

}

\newsitem{火星探测器进入安全模式：冗余系统并非无限保险}{重点报道｜Space.com｜11月21日}{
火星上的毅力号漫游车和环绕火星运行的奥德赛号近期分别遇到通信与姿态控制问题，科学观测受到影响。报道指出，毅力号主用超高频无线电出现异常后转用备用设备，并使用直接对地通信进行诊断；奥德赛号则因反作用轮故障进入安全模式。两种故障不同，却共同说明长期无人探测任务对备份系统的依赖。\cite{mars_fault}

漫游车通常借助火星轨道器中继数据，因为中继链路能够在合适的几何条件下提供更有效的通信。直接对地通信虽可作为退路，但速率和时机受到限制。轨道器使用反作用轮调整姿态，可以在不频繁消耗推进剂的情况下控制指向；若反作用轮失效，改用推进器可能显著缩短剩余任务寿命。

安全模式的意义是优先保护设备，而不是宣告任务失败。工程师需要先确认供电、热控和通信稳定，再逐步恢复观测。对公众而言，判断任务是否真正恢复，应看NASA的正式状态更新，而不是仅根据团队“希望月底恢复”的目标时间。

}

\newsitem{数学拾趣：本周预印本筛查与证明可信度}{方法栏目｜arXiv 数学、Lean/mathlib｜11月15—21日}{
本周对数论、代数、几何、组合、逻辑、概率和形式化证明方向进行了检索，但未能获得足够可靠、可逐项核对题名、作者、版本和具体定理陈述的当周预印本，因此不刊登未经核实的所谓重大数学突破。这不意味着本周没有值得阅读的论文，而是本次筛查证据尚不足以达到正式发表标准。

理解数学新闻，需要区分三个层次：作者在预印本中声称证明某个命题；其他数学家对论证进行独立检查；最终形成经过修订、可复现且被共同体接受的证明。即使出现大量Lean代码，也需要检查形式化语句是否与原命题一致、是否依赖未经证明的公理、编译是否可重复，以及形式化对象是否包含关键数学假设。

例如，若某人声称证明黎曼猜想，单凭一个能够编译的Lean文件不能说明结论成立；必须明确它究竟证明了复平面上黎曼ζ函数所有非平凡零点实部等于二分之一，还是证明了一个条件性、有限范围或被弱化的命题。数学形式化能减少证明步骤中的逻辑疏漏，却不能自动替代命题建模和数学解释。

本栏目将继续优先寻找有清晰定义、可计算例子、实质方法创新和可追溯源文件的成果。没有足够证据时保留空缺，比制造轰动标题更符合数学研究的习惯。

}

\end{documents}
\color{black}
\begingroup\footnotesize
\bibliographystyle{plain}
\bibliography{refs}
\endgroup
\end{document}
